\begin{helper}
Let \(C\) be a nonempty finite set of candidates.  Fix an activation
parameter \(p\in[0,1]\) and candidate priorities \(\pi(x)\in[0,1]\) for
\(x\in C\).  Let \(\eta\) be the product law of one Bernoulli-\(p\)
activation coordinate together with one independent uniform priority
coordinate, with total mass \(1\), expressed by the lower-rectangle identities
\[
\eta\{a=1,\ 0\le q\le t\}=\operatorname{ofReal}(pt)
\qquad(0\le t\le1)
\]
and \(\eta\{a=1,\ q\in[0,1]\}=\operatorname{ofReal}(p)\).
Let \(\theta\) be the probability law of independent Bernoulli-\(p\)
activation coordinates and independent uniform priority coordinates indexed by
\(C\):
for every \(A\subseteq C\) and every \(t_x\in[0,1]\), \(x\in A\),
\[
\theta\{a_x={\bf 1}_{x\in A}\text{ for all }x\in C,\ 
       0\le q_x\le t_x\text{ for all }x\in A\}
=\operatorname{ofReal}\!\left(
  p^{|A|}(1-p)^{|C|-|A|}\prod_{x\in A}t_x\right).
\]
Then the probability that at least one candidate is not killed by its private
blocker under \(\theta\) is at least the probability that at least one
candidate is not killed by the single common blocker under \(\eta\):
\[
\eta\{\neg(a=1,\ q\in[0,1],\ \pi(x)<q\text{ for every }x\in C)\}
\le
\theta\{\exists x\in C:\neg(a_x=1,\ q_x\in[0,1],\ \pi(x)<q_x)\}.
\]
\end{helper}
\begin{proof}
Apply
\(\noderef{SamplingCandidateSetReplacementSurvivalComparison}\) to the finite
set \(C\), the number \(p\), and the function \(\pi\).  Its statement supplies
a witness \(x_0\in C\) such that \(\pi(x)\le \pi(x_0)\) for every \(x\in C\)
and
\[
1-p(1-\pi(x_0))
\le
1-\prod_{x\in C}p(1-\pi(x)).
\tag{1}
\]
We use this supplied witness in the one-blocker computation; no separate
maximizer is chosen.

Let
\[
K=\{a=1,\ q\in[0,1],\ \pi(x)<q\text{ for every }x\in C\}
\]
be the event that the common blocker kills all candidates.  Because \(x_0\)
has maximal priority, the condition \(\pi(x)<q\) for every \(x\in C\) is
equivalent to \(\pi(x_0)<q\).  Inside the active unit event
\[
U=\{a=1,\ 0\le q\le1\},
\]
the two disjoint events
\[
L=\{a=1,\ 0\le q\le\pi(x_0)\}
\quad\text{and}\quad
K=\{a=1,\ \pi(x_0)<q\le1\}
\]
partition \(U\).  This is an exact order partition; it does not use
conditioning on the value \(q=\pi(x_0)\) and does not require a separate
atomlessness assertion.  The hypotheses on \(\eta\) give
\[
\eta(U)=\operatorname{ofReal}(p),\qquad
\eta(L)=\operatorname{ofReal}(p\,\pi(x_0)),
\]
since \(0\le\pi(x_0)\le1\).  As \(0\le p\,\pi(x_0)\le p\le1\), finite
additivity on the disjoint union \(U=L\sqcup K\) gives
\[
\eta(K)
=\operatorname{ofReal}(p)-\operatorname{ofReal}(p\,\pi(x_0))
=\operatorname{ofReal}\bigl(p(1-\pi(x_0))\bigr).
\]
The common-blocker survival event in the statement is \(K^c\).  Since
\(\eta\) has total mass \(1\),
\[
\eta(K^c)=\operatorname{ofReal}\bigl(1-p(1-\pi(x_0))\bigr).
\tag{2}
\]

Now let \(P\) be the event that every private blocker kills its candidate:
\[
P=\{a_x=1,\ q_x\in[0,1],\ \pi(x)<q_x
       \text{ for every }x\in C\}.
\]
The private-blocker survival event in the statement is \(P^c\).  Indeed, over
the finite nonempty set \(C\), saying that some candidate is not killed by its
private blocker is exactly the negation of saying that all candidates are
killed by their private blockers.

It remains to compute \(\theta(P)\) from the cylinder identities.  Write
\[
U_C=\{a_x=1,\ 0\le q_x\le1\text{ for every }x\in C\}
\]
for the all-active unit cube.  Taking the activation set \(A=C\) and the
constant threshold \(t_x=1\) in the cylinder identity gives
\[
\theta(U_C)=\operatorname{ofReal}(p^{|C|}).
\]
Apply \noderef{SamplingFiniteProductLowerOrthantEqualityMeasurable}
with \(F=U_C\), the priority projection \(q=(q_x)_{x\in C}\),
\(M=\operatorname{ofReal}(p^{|C|})\), and the box
\(D=\prod_{x\in C}(\pi(x),1]\).
The projection is measurable because the supplied sample space has the
all-subsets measurable structure.  The preceding threshold-one calculation
gives \(\theta(F)=M<\infty\).  For unit thresholds \(t_x\), the condition
\(0\le q_x\le t_x\) already implies \(q_x\in[0,1]\); hence the cylinder
identity with \(A=C\) gives exactly
\[
\theta\{F,\ 0\le q_x\le t_x\text{ for all }x\}
=M\operatorname{ofReal}\left(\prod_{x\in C}t_x\right).
\]
The box \(D\) is Borel and lies in the unit cube since \(\pi(x)\ge0\).
The cited theorem therefore gives the mass of \(F\cap\{q\in D\}\),
with its redundant unit-cube intersection, as \(M\operatorname{vol}(D)\).
This event is exactly \(P\).  Finite-product Lebesgue volume gives
\(\operatorname{vol}(D)=\prod_{x\in C}\operatorname{ofReal}(1-\pi(x))\).
All factors are nonnegative, so their product is
\(\operatorname{ofReal}(\prod_{x\in C}(1-\pi(x)))\).  Since
\(\prod_{x\in C}p(1-\pi(x))=p^{|C|}\prod_{x\in C}(1-\pi(x))\), we get
\[
\theta(P)=
\operatorname{ofReal}\!\left(\prod_{x\in C}p(1-\pi(x))\right).
\tag{4}
\]
All terms are nonnegative and at most \(1\), because \(0\le p\le1\) and
\(0\le\pi(x)\le1\).  Therefore total mass \(1\) and (4) give
\[
\theta(P^c)=
\operatorname{ofReal}\!\left(
1-\prod_{x\in C}p(1-\pi(x))\right).
\tag{5}
\]

Combining (1), (2), and (5), and using monotonicity of
\(\operatorname{ofReal}\) on real numbers, yields exactly
\[
\eta(K^c)\le\theta(P^c).
\]
This is the claimed comparison.  The argument only uses finite additivity,
the stated lower-rectangle/cylinder identities, and the cited measurable
lower-orthant theorem; it never conditions on an exact value of a continuous
priority coordinate and never infers arbitrary strict-threshold event
probabilities directly from \(\noderef{RandomIndependentSetSampling}\).
\end{proof}
